% Loopotron: depth-wise weight sharing in a hybrid Mamba-Transformer MoE stack.
%
% Compile with: pdflatex loopotron.tex  (twice, for references)
% Dependencies are deliberately limited to standard TeX Live packages.
%
% NOTE: this file has not been compiled -- no LaTeX toolchain was available on the machine where
% it was written. If a package is missing, the likely culprits are booktabs and hyperref.

\documentclass[11pt,a4paper]{article}

\usepackage[margin=2.4cm]{geometry}
\usepackage{booktabs}
\usepackage{amsmath}
\usepackage{graphicx}
\usepackage{xcolor}
\usepackage[hidelinks]{hyperref}

% Layer-pattern strings contain ^, * and brackets; \lp{} keeps them typeset literally.
\newcommand{\lp}[1]{\texttt{#1}}
% \textasciicircum, not \^{}: the latter is a circumflex *accent* and renders raised and undersized
% in typewriter font, which matters because the layer-pattern strings are load-bearing here.
\newcommand{\caret}{\textasciicircum}

\title{\textbf{Loopotron}\\[0.3em]
\large Where is extra depth worth spending in a hybrid Mamba--Transformer MoE?}
\author{}
\date{Experiments run 2026-08-06 -- 2026-08-11}

\begin{document}
\maketitle

\begin{abstract}
\noindent
A Nemotron-style hybrid language model is a stack of single-operator residual layers --- Mamba-2
mixers, sparse mixture-of-experts feed-forwards, and grouped-query attention --- described by a
pattern string. Because each layer is a self-contained residual block, a \emph{loop} (executing a
run of layers several times while reusing one set of weights) can be placed around any subset of
layer types. This trades parameters for effective depth, and lets us ask where extra depth is worth
spending. We train ten models at a matched budget of 557M tokens: an unlooped hybrid baseline,
three looped variants, three conditioning refinements, an iso-FLOP fresh-weight anchor, and two
vanilla dense transformers. Three results survive a measured seed-noise floor. First, \textbf{every
looped variant beats the unlooped hybrid} on held-out loss, by 12--48 standard deviations of
run-to-run noise. Second, \textbf{the fresh-weight anchor still beats every looped variant} at
matched FLOPs, reproducing the published perplexity penalty for weight sharing. Third, and
largest, \textbf{the hybrid beats an iso-FLOP dense transformer by 0.198 nats}. We also report two
negative results: per-iteration normalization and input injection have no measurable effect, and a
synthetic reasoning suite designed to detect the reasoning-over-perplexity trade-off has
insufficient resolution at this budget --- a claim we initially made from it, and here retract.
\end{abstract}

\section{Introduction}

Nemotron-style hybrid models~\cite{nemotron} interleave three operator types according to a layer
pattern string, in which \lp{M} denotes a Mamba-2 mixer, \lp{E} a sparse mixture-of-experts
feed-forward, \lp{*} grouped-query attention, and \lp{-} a dense feed-forward. Unlike a classical
transformer block, which bundles attention and a feed-forward network, every layer here wraps
exactly one operator in a pre-norm residual connection,
\begin{equation}
x \leftarrow x + f(\mathrm{norm}(x)),
\end{equation}
so the ratio of Mamba to attention to MoE layers is tunable independently.

This structure makes a natural experiment available. Because each layer is self-contained, a
\emph{loop} --- executing a run of layers several times while reusing one set of weights, as in a
Universal Transformer --- can be placed around any subset of layer types. Looping holds parameters
fixed while raising FLOPs, so it converts parameters into effective depth. The question this work
asks is: \textbf{where is that depth worth spending --- on the Mamba mixers, the MoE
feed-forwards, the attention layers, or on mixed blocks?}

The published consensus is that depth-wise weight sharing loses on perplexity and wins on
reasoning~\cite{saunshi,lssm}. That expectation shaped the design: we measure held-out loss and a
synthetic reasoning suite together, and we include a fresh-weight \emph{anchor} for the leading
looped variant so that ``more compute helps'' can be separated from ``sharing works''.

\section{Method}

\subsection{Loop grammar}

A loop group is written \lp{[\textless symbols\textgreater]\caret K} inside a layer pattern; for
example \lp{M[ME]\caret3*E} builds one Mamba and one MoE layer and executes the pair three times.
The repeat operator is \lp{\caret} rather than \lp{*}, because \lp{*} is the attention symbol and
\lp{M[ME]*3*E} would read as though the 3 were multiplied by an attention layer. Nesting is
unsupported. Bracket-free patterns parse exactly as before, so all pre-existing configurations are
unaffected. Throughout, \emph{built} layers count sets of weights and \emph{executed} layers count
applications per token, so \lp{M[ME]\caret3E} has 4 built and 8 executed layers.

The implementation keeps the layer container flat, with one entry per set of weights, and puts the
loop in a separate execution schedule. This is what allows activation checkpointing, FSDP2 block
wrapping, the weight-decay group regular expressions and the initialization parameter filters to
work on looped models unchanged, since parameter names are untouched.

Three quantities must follow \emph{executions} rather than parameters, and getting any of them
wrong biases an ablation silently rather than raising an error:
\begin{itemize}
  \item \textbf{Effective depth} for depth-scaled initialization. Residual variance grows with the
    number of layer applications, so a looped variant must be initialized for its executed depth.
    This is the single most important detail; the configuration derives it from the pattern so it
    cannot go stale.
  \item \textbf{The MoE auxiliary load-balancing loss}, which is accumulated across visits and
    reset per forward pass. A looped MoE layer routes independently on each visit.
  \item \textbf{Active parameters} for the FLOPs estimate, which are execution-weighted.
\end{itemize}

\subsection{The research model}

All models are a geometric scale-down of Nemotron-3 Nano 30B-A3B that preserves the reference's
\emph{cost structure}, which is what makes conclusions about where to spend depth transferable
(Table~\ref{tab:scaledown}).

\begin{table}[h]
\centering
\begin{tabular}{lrr}
\toprule
Property & Reference 30B-A3B & This work \\
\midrule
Model dimension / built layers      & 2688 / 52          & 1024 / 12 \\
Parameters                          & 31.58B / 3.58B act.& 1.105B / 225M act. \\
Non-embedding active / total        & 9.3\%              & 9.7\% \\
Per-layer active cost $M : * : E$   & $1 : 0.60 : 2.07$  & $1 : 0.59 : 2.01$ \\
Expert sparsity                     & top-6 of 128       & top-6 of 128 \\
\bottomrule
\end{tabular}
\caption{The research model preserves the reference's per-layer cost ratios and expert sparsity.}
\label{tab:scaledown}
\end{table}

The model is 1.1B rather than 400M by necessity: the embedding table and the Mamba layers are dense
and always active, so a 400M model with a 128k vocabulary cannot be $\sim$10\% active, and forcing
the size down distorts the $M : * : E$ cost ratio, which is exactly what the ablation is about.
Training FLOPs scale with \emph{active} parameters, so these models train at roughly the speed of a
225M dense model; only memory holds 1.1B.

\subsection{Training setup}

Every run is identical except where stated: FineWeb (2B-token shard), sequence length 2048, global
batch $8 \times 4 \times 2048 = 65{,}536$ tokens per step, AdamW ($\beta = (0.9, 0.95)$, weight
decay 0.1 excluding embeddings, norms, SSM parameters and the router), linear warmup over 750 steps
followed by cosine annealing from $4.5\times10^{-4}$ to $4.5\times10^{-5}$, bf16 mixed precision,
full activation checkpointing, one A100-80GB per run. All models train for
\textbf{8500 steps = 557M tokens}, and the learning-rate schedule completes within that budget.

That last point is not a detail. An earlier wave derived the schedule length from the size of the
training file (30{,}517 steps) while stopping runs on wall clock at 5--11k steps, so the cosine
never decayed: the learning rate was still at 96\% of its peak at step 4750 and the loss flattened
at 3.63 against the 3.34 reached once the schedule completed. Worse for an ablation than for the
absolute number, a constant peak learning rate puts every variant on the same noisy plateau. All
results below use the completed schedule.

\subsection{Evaluations}

Because the published effect is expected to live in reasoning rather than perplexity, we log four
held-out measures alongside the language-modelling loss. Two are real benchmarks scored by the
negative log-likelihood of the answer span: \textbf{Minerva MATH} (4-shot) and closed-book
\textbf{TriviaQA}, the latter acting as a parametric-knowledge control, since closed-book recall
lives in weights that a loop does not add. Accuracy on both is meaningless at this scale and is not
reported.

Two are synthetic and deterministic, parameterized by a \emph{hop count} --- the number of serial
lookups the answer requires --- over a 26-symbol alphabet, giving a chance accuracy of
$1/26 = 3.85\%$. In \textbf{$p$-hop induction}, a sequence of 256 random symbols is presented; the
last symbol is the query, and one hop finds its most recent earlier occurrence and reads off the
symbol that follows it. Each further hop repeats this, searching only to the left of the previous
match, so hop $k+1$ cannot begin until hop $k$ has produced both its symbol and its position. One
hop is exactly the induction-head operation and acts as the control. For example, over a 6-symbol
alphabet:
\begin{center}
\begin{tabular}{@{}l@{}}
\ttfamily E~A~B~B~B~E~F~D~A~A~B~C~D~C~B~A~E~E~A~A \\[-0.2em]
\ttfamily \phantom{E~A~B~B~B~E~F~D~A~A~B~C~D~C~B~}$\uparrow$\phantom{A~E~}$\uparrow$~$\uparrow$ \\[-0.2em]
\ttfamily \phantom{E~A~B~B~B~E~F~D~A~A~B~C~D~C~B~}\,2\phantom{A~E~}\,1~q \\
\end{tabular}
\end{center}
Reading right to left: \texttt{q} marks the query \texttt{A}; hop 1 finds its most recent earlier
occurrence and takes the symbol following it, again \texttt{A}; hop 2 searches only further left,
finds \texttt{A} at position 15 and takes \texttt{E}, the 2-hop answer. In \textbf{variable binding}, a
\emph{shuffled} list of assignments dereferences a chain to one literal, so recency heuristics do
not help; a real sample is
\begin{center}
\ttfamily ; F = S ; E = H ; L = A ; T = U ; P = G ; Z = V ; N = D ; H = N ; O = K ; W = Y ; J = C ; Q = R ; E =
\end{center}
whose answer is \texttt{D} via $E \to H \to N \to D$. This task has a format-aware shortcut: the
answer is always a symbol appearing as a value but never as a variable, and there are exactly ten
such, so a guesser exploiting the format reaches 10\% without following the chain. It produced no
signal in any run and is not reported below.

Every sample is emitted pre-shifted and masked so that only the answer position is scored; the
reported figure is therefore the negative log-likelihood, or accuracy, of the answer alone, with
each scored token weighted equally rather than each batch.

\subsection{Variants}

Table~\ref{tab:main} lists the models. The looped variants are matched on \textbf{FLOPs, not loop
count}: an extra MoE execution costs about twice an extra Mamba execution and attention about
1.65$\times$, so equal $K$ across types would compare unequal compute budgets. The loop is also
spread over \emph{every} layer of the relevant type rather than concentrated on one, so ``which
layer happened to be looped'' is not a confound.

Two comparisons matter. Against the unlooped baseline \textbf{A0} (iso-parameter) the question is
whether looping buys usable compute. Against the \emph{anchor} \textbf{N4} --- a bracket-free model
with the same executed count of every layer type, the same active parameters and the same FLOPs as
A6, differing only in that its executions use fresh weights --- the question is whether sharing is
as good as not sharing. Two vanilla dense transformers (pre-norm, RoPE, SwiGLU, multi-head
attention, tied embeddings) supply the reference the hybrid literature usually omits: \textbf{D1}
matches A0's \emph{active} parameters and therefore its FLOPs per token, while \textbf{D2} matches
A0's \emph{total} parameters and consequently spends about 4.9$\times$ the compute per token.

Finally, three conditioning refinements are applied to A6. \textbf{Per-iteration normalization}
gives each loop iteration its own pre-norm while the operator stays shared, adding $(K-1) \cdot
d_{\text{model}}$ parameters per looped layer (4096 in total here). \textbf{Input injection} adds
the loop group's input back to the hidden states at the start of every iteration after the first,
adding no parameters. Both are described in the literature as near-prerequisites for scaling past
two or three loop iterations, and both are deliberately optional so that weight sharing and
per-iteration conditioning remain separable.

\section{Results}

\subsection{Main comparison}

\begin{table}[t]
\centering
\small
\setlength{\tabcolsep}{4pt}
\begin{tabular}{llcrrcccc}
\toprule
& & Built/ & & & & \multicolumn{3}{c}{Held-out} \\
\cmidrule(l){7-9}
Model & Layer pattern & exec. & Params & Active & FLOPs & LM loss & MATH & TriviaQA \\
\midrule
\multicolumn{9}{l}{\emph{Dense transformer baselines}} \\
D1 & 9 layers, $d{=}1024$   & 9/9   & 226M   & 226M   & 1.00 & 3.5383 & 4.4253 & 6.7404 \\
D2 & 21 layers, $d{=}2048$  & 21/21 & 1.111B & 1.111B & 4.90 & \textbf{3.2591} & \textbf{4.1053} & \textbf{6.2516} \\
\midrule
\multicolumn{9}{l}{\emph{Hybrid, no loop}} \\
A0 & \lp{MEM*EMEMEM*E}      & 12/12 & 1.105B & 226M & 1.000 & 3.3404 & 4.1415 & 6.4363 \\
N4 & \lp{MEM*E*EMEMEM*E*E}  & 16/16 & 1.487B & 256M & 1.179 & \textbf{3.3021} & 4.0818 & 6.2288 \\
\midrule
\multicolumn{9}{l}{\emph{Looped variants} (ordered by LM loss)} \\
A5 & \lp{ME[M*]\caret3EMEME[M*]\caret3E}            & 12/20 & 1.105B & 262M & 1.261 & 3.3168 & 4.0949 & 6.4924 \\
A4 & \lp{[ME]\caret2M*E[ME]\caret2[ME]\caret2M*E}   & 12/18 & 1.105B & 278M & 1.220 & 3.3222 & 4.1754 & 6.4295 \\
A6 & \lp{MEM[*E]\caret2MEMEM[*E]\caret2}            & 12/16 & 1.105B & 256M & 1.179 & 3.3243 & 4.1569 & 6.4342 \\
A6a & \quad + per-iteration norm                    & 12/16 & 1.105B & 256M & 1.179 & 3.3246 & 4.1681 & 6.3929 \\
    & \multicolumn{5}{r}{\emph{(4 seeds, s.d.)}}                                    & $\pm$0.0005 & $\pm$0.0177 & $\pm$0.0629 \\
A6c & \quad + norm and injection                    & 12/16 & 1.105B & 256M & 1.179 & 3.3301 & 4.2102 & 6.4325 \\
A6b & \quad + input injection                       & 12/16 & 1.105B & 256M & 1.179 & 3.3339 & 4.2258 & 6.5063 \\
\bottomrule
\end{tabular}
\caption{All models at 8500 steps (557M tokens), identical data, schedule and evaluations.
``Built/exec.'' is sets of weights over layer applications per token; FLOPs are relative to A0 per
token. LM loss is held-out cross-entropy; MATH and TriviaQA are the negative log-likelihood of the
answer span (lower is better throughout). D2 is not an iso-FLOP comparison. A6a and A6c carry 4096
extra parameters for their per-iteration norms.}
\label{tab:main}
\end{table}

\subsection{The seed-noise floor}

Every number in Table~\ref{tab:main} except A6a's is a single run, so the differences are
meaningless without a measure of run-to-run variation. Four runs of A6a from identical
configurations --- the weight initialization is unseeded, so each run draws different weights while
the data order is held fixed --- give Table~\ref{tab:noise}.

\begin{table}[h]
\centering
\begin{tabular}{lrrr}
\toprule
Metric & Mean & s.d. & Range \\
\midrule
Held-out LM loss        & 3.3246 & \textbf{0.0005} & 0.0011 \\
$p$-hop-1 accuracy      & 0.0952 & \textbf{0.0172} & 0.0708 -- 0.1079 \\
$p$-hop-1 nll           & 3.276  & 0.114  & 0.245 \\
Minerva MATH nll        & 4.1681 & 0.0177 & --- \\
TriviaQA nll            & 6.3929 & 0.0629 & --- \\
\bottomrule
\end{tabular}
\caption{Four replicates of A6a. The two instruments differ in resolution by more than two orders
of magnitude relative to the effects being measured.}
\label{tab:noise}
\end{table}

The consequences point in opposite directions.

\paragraph{Held-out loss resolves the ordering.} At a standard deviation of $0.0005$, every
difference in Table~\ref{tab:main} is far outside noise. N4 versus A0 (0.038) is roughly 76 standard
deviations; even the smallest gap, A6b versus A0 (0.006), is roughly 12. The ordering
anchor $>$ looped variants $>$ unlooped baseline $>$ iso-FLOP dense is therefore real.

\paragraph{The reasoning suite does not resolve anything.} The seed standard deviation on
$p$-hop-1 accuracy (0.0172) exceeds the entire spread across hybrid variants, and the standard
deviations on MATH ($\pm0.0177$) and TriviaQA ($\pm0.0629$) cover every difference between hybrid
rows in Table~\ref{tab:main}. Resolving a 0.02 difference in $p$-hop accuracy at two standard
deviations would require roughly twelve runs per variant. Crucially, more evaluation questions
would not help: the sampling standard error over 2048 questions is 0.007, so the dominant variance
is in the model, not the question sample.

\subsection{Findings}

\paragraph{1. Looping helps, modestly and reliably.} Every looped variant beats the unlooped
baseline A0 by 0.006--0.024 nats, or 12--48 standard deviations. The best is A5, which loops
Mamba-plus-attention pairs three times.

\paragraph{2. Fresh weights still beat shared weights at matched FLOPs.} N4 (3.3021) beats A6
(3.3243), its exact iso-FLOP anchor, by 0.022 nats. This reproduces the published perplexity
penalty for depth-wise weight sharing. What the loop buys instead is memory: A6 reaches the same
executed depth with 1.105B parameters against N4's 1.487B.

\paragraph{3. The hybrid beats an iso-FLOP dense transformer by 0.198 nats}, comparing A0 to D1.
This is by far the largest effect measured, and it is the comparison most often missing from
hybrid-architecture ablations. It comes with a caveat: D1 trains at 34.5 samples/s against A0's
23.3 (MFU 0.36 versus 0.22), because dense attention and a plain SwiGLU use the hardware far better
than Mamba and MoE kernels. Per unit of \emph{wall clock} rather than of FLOPs, the hybrid's
advantage is substantially smaller.

\paragraph{4. Per-iteration normalization and input injection do nothing measurable.} A6a
(3.3246), A6c (3.3301) and A6b (3.3339) straddle A6 (3.3243). We initially reported that
per-iteration norms improved $p$-hop-1 accuracy from 0.073 to 0.108, a nominal 3.9 standard
deviations. That figure used the \emph{sampling} standard error over the evaluation set (0.007) in
place of run-to-run noise (0.017), which it understates by a factor of 2.5. Under replication,
0.108 proved to be the top of A6a's own seed range and A6's 0.073 essentially the bottom of it; the
difference is 1.2 standard deviations. \textbf{We retract the claim.} Sampling error is a lower
bound on measurement error, never a substitute for it.

\paragraph{5. The expected reasoning-over-perplexity dissociation did not appear.} MATH and
TriviaQA were included expecting them to move in opposite directions if a variant bought reasoning
without buying parameters. They do not: both track held-out loss almost perfectly across all ten
models, with N4 among the best and D1 the worst on all three. No variant is trading knowledge for
reasoning at this budget.

\paragraph{6. Block-loop underperforms layer- and pair-loop.} A4, which repeats a Mamba-plus-MoE
block, trails on every reasoning measure, consistent with reports that repeating each layer in
place beats repeating the stack.

\paragraph{7. A dense-model anomaly, resolved by scale.} D1 scores 0.0073 average $p$-hop accuracy
--- reliably \emph{below} the 0.0385 chance level, which requires structure rather than ignorance
--- while scoring normally on variable binding, which uses the same alphabet and the same
evaluation path. D2, the same architecture with 4.9$\times$ the compute, does not collapse
(0.0527). The effect is therefore one of scale and budget, not of architecture class: a small dense
transformer at 557M tokens has not formed the in-context machinery to notice that the answer must
be one of the 26 symbols present, whereas Mamba's recurrent state tracks the in-context
distribution essentially for free, and every hybrid lands near $\ln 26 = 3.26$ nats.

\section{Limitations}

\paragraph{Budget.} 557M tokens is roughly 12\% of a compute-optimal allocation for a 225M-active
model. Differences of 0.01--0.04 nats at this scale need not survive to convergence, and the
published loop results are reported at 250B tokens and above.

\paragraph{Seeds.} Only A6a is replicated, and its four runs vary weight initialization only, with
the data order held fixed. Run-to-run noise under a varying data seed can only be larger, so
Table~\ref{tab:noise} is a lower bound. Every other row of Table~\ref{tab:main} is $n=1$.

\paragraph{Incomplete arm coverage.} The three single-layer-type loops --- Mamba alone at $K=3$,
MoE alone at $K=2$, attention alone at $K=4$ --- have not yet been run with the evaluation suite,
and they are precisely the variants that would answer the titular question most directly. Only
mixed-block loops appear above. Anchors for A4 and A5 have also not been constructed.

\paragraph{Reasoning measures.} The synthetic suite is under-powered here, and the $p$-hop task may
additionally be mis-scaled: at a prompt length of 256 over a 26-symbol alphabet the query appears
about 9.8 times, so even a perfectly functioning induction head returns a mixture over roughly ten
candidate successors and lands near 10\% --- which is approximately where the best models sit. A
shorter prompt would make the most-recent match near-unique and give the task room to discriminate.

\section{Conclusion}

At a matched 557M-token budget, depth-wise weight sharing in a hybrid Mamba--Transformer MoE stack
gives a small but reliable improvement in held-out loss over an unlooped model of the same size,
while remaining behind an iso-FLOP model with fresh weights --- buying memory rather than quality.
The hybrid architecture itself, however, is worth considerably more than the loop: it beats an
iso-FLOP dense transformer by an order of magnitude larger a margin than any difference among the
looped variants, though part of that advantage is returned as lower hardware utilization.

The methodological finding may be the more useful one. Held-out loss and the reasoning suite differ
in measurement resolution by more than two orders of magnitude relative to the effects of interest,
and a nominally 3.9-sigma result computed from sampling error alone did not survive four
replications. Any claim from a suite of this kind should be accompanied by a measured seed-noise
floor before it is believed --- including, and especially, one's own.

\begin{thebibliography}{9}
\bibitem{nemotron} NVIDIA. \emph{Nemotron-3 Nano}. arXiv:2512.20848.
\bibitem{saunshi} Saunshi et al. \emph{Reasoning with Latent Thoughts: On the Power of Looped
  Transformers}. arXiv:2502.17416.
\bibitem{lssm} \emph{Looped State-Space Language Models}. arXiv:2607.10110.
\bibitem{phi} \emph{How Much Is One Recurrence Worth?} arXiv:2604.21106.
\bibitem{loopie} \emph{Loopie: Looped Mixture-of-Experts}. arXiv:2607.16051.
\end{thebibliography}

\end{document}
